% Amendment A3, source: results/phase2/a3/analysis/final/{table_a3.md,a3_paired.csv}
\textbf{Repairing the OOD gate (pre-registered amendment A3).} The H2 failure suggested a diagnosis: the raw detector
scores every input change, including sensor noise the policy has learned to ignore, and it does not consider which action
is proposed. We therefore replaced only the OOD score with an \emph{action-conditioned} one. It is a class-conditional
Mahalanobis distance in the style of~\cite{lee2018mahalanobis}, computed on each member's last hidden layer (128-d), with
the ensemble's proposed action as the class and the score averaged over members. It uses the same shrinkage, the same
10-step window, the same fit maps and the same ID-validation 95th-percentile threshold. Design, decision rule and
comparisons were committed before the detector was fitted. Nothing was tuned, and the amended code again reproduced a
frozen cell exactly. Under sensor shift the new gate fires on 5.4\% of decisions instead of 93.9\%, and under combined
shift on 20.3\% instead of 95.5\% (Fig.~\ref{fig:a3}). Relative to the raw gate, success under sensor shift rises by 3.3
points (0.927$\to$0.959). Attempted unsafe moves fall by 0.28 (sensor) and 0.24 (combined) per episode, and collisions
fall by 0.39 and 0.66. No primary metric becomes worse in any condition, so H2$'$ is supported. The second pre-registered
comparison, however, finds no added value: against the same system with \emph{no} OOD gate, success and unsafe moves do
not differ significantly in any condition, and the only significant differences are 0.7--3.2 points more interventions
and 0.2--0.3 units less energy. Used alone (gate + fallback, no monitor or uncertainty gate), the action-conditioned
detector rarely fires (1.9--9.9\% of decisions) and leaves 1.0--2.0 attempted unsafe moves per episode. The repair
therefore removes the harm, but it does not make OOD detection a useful safety signal here.

\begin{figure}[t]
\centering
\includegraphics[width=\columnwidth]{fig9_a3}
\caption{Amendment A3: paired differences of the full system with the action-conditioned OOD gate against the raw OOD
gate (blue) and against no OOD gate (orange). Mean, 95\% CI over 10 seeds; * Holm $p<0.05$.}
\label{fig:a3}
\end{figure}
